\documentclass[12pt]{article}
\usepackage{amsmath,amssymb,booktabs,longtable,array}
\usepackage{fontspec}
\usepackage{polyglossia}
\setmainlanguage{arabic}
\setotherlanguage{english}
\newfontfamily\arabicfont[Script=Arabic]{Amiri}
\usepackage{geometry}
\geometry{margin=2.1cm}
\usepackage{hyperref}

\title{مقارنة تجارب المشروع مباشرة بتنبؤات Leung 2026 وHardy--Littlewood}
\author{مشروع تجربة البحث عن الأعداد الأولية}
\date{2026-09-01}

\begin{document}
\maketitle

\section{الهدف}

قبل الانتقال من الأزواج إلى نموذج ثلاثي \(k=3\)، نختبر سؤالًا منهجيًا:

\begin{quote}
هل السلوك الذي رأيناه بالفعل يحتاج إلى معلومات ثلاثية، أم أن نظرية الارتباط الزوجي والتقريب الغاوسي متعدد المتغيرات تفسره قبل أن نضيف بنية جديدة؟
\end{quote}

لذلك نفذنا تجربتين:

\begin{enumerate}
\item تجربة مصطفة مع صياغة Leung 2026 باستعمال العد الموزون
\(\psi\) في فترات قصيرة ذات طول نسبي ثابت.
\item إعادة فحص سلسلة الفترات بين مربعات الأوليات باستعمال نموذج Gaussian
تحدده مصفوفة الارتباط الزوجي فقط، ثم اختبار runs والأنماط الثلاثية.
\end{enumerate}

\section{تنبؤ Leung 2026}

يدرس Leung العد الموزون
\[
\psi(x)=\sum_{p^k\le x}\log p.
\]

لفترات طولها
\[
h=\delta x,
\]
يثبت، تحت RH وLI، توزيعًا Gaussian متعدد المتغيرات ذا ارتباطات سالبة ضعيفة.

لفترتين متجاورتين يعطي Corollary 1.1، في الحد الرئيسي:
\[
\rho_1
\approx
-\frac{\log 2}{\log(1/\delta)}.
\]

وتعطي Proposition 2.2 صيغة أدق للتغاير:
\[
\mathrm{Cov}_{jk}
=
-\Delta(|t_j-t_k|)\delta
+
O\!\left((T\delta\log(1/\delta))^2\right),
\qquad j\ne k,
\]
حيث
\[
\Delta(t)
=
\frac12\left(
(t+1)\log(t+1)
-2t\log t
+(t-1)\log(t-1)
\right).
\]

أما التباين فهو، إلى الحد الثانوي:
\[
V
=
\delta\log\frac1\delta
+
(1-\gamma-\log2\pi)\delta
+\cdots.
\]

وبالتالي استعملنا للمقارنة المنتهية:
\[
\rho_k^{(\mathrm{sec})}
=
-\frac{\Delta(k)}
{\log(1/\delta)+1-\gamma-\log2\pi}.
\]

هذه الصيغة الأخيرة ليست حدًا نهائيًا جديدًا؛ هي مجرد استعمال مباشر للحدين
الرئيسيين في صيغة Leung لتقليل انحياز الحجم المنتهي.

\section{التجربة المصفوفة مع Leung}

اخترنا \(8000\) نقطة \(x\) موزعة بالتساوي في \(\log x\) داخل المجال:
\[
20\,000\,000\le x\le55\,000\,000.
\]

لكل:
\[
\delta\in\{0.001,0.002,0.004,0.008\},
\]
حسبنا أربع فترات متجاورة طول كل منها \(\delta x\)، واستعملنا العد الموزون
\(\psi\)، ثم قسمنا الانحراف على التباين المتوقع من صيغة Leung.

\subsection{التباين والشكل الهامشي}

كان متوسط تباين المتغير المعياري عبر الفترات الأربع:

\begin{center}
\begin{tabular}{rrr}
\toprule
\(\delta\) & متوسط التباين المرصود & متوسط kurtosis\\
\midrule
0.001 & 1.0724 & 3.0252\\
0.002 & 1.0544 & 2.8317\\
0.004 & 1.0208 & 2.8066\\
0.008 & 0.9824 & 3.1176\\
\bottomrule
\end{tabular}
\end{center}

إذن التطبيع النظري يضع التباين قريبًا من \(1\)، والشكل الهامشي قريبًا
من Gaussian على هذا المدى، مع انحرافات منتهية متوقعة.

\subsection{الارتباط المتجاور}

\begin{center}
\begin{tabular}{rrrr}
\toprule
\(\delta\) & \(\rho_1\) المرصود &
تنبؤ الحد الرئيسي & تنبؤ الحد الثانوي\\
\midrule
0.001 & -0.1410 & -0.1003 & -0.1262\\
0.002 & -0.1696 & -0.1115 & -0.1444\\
0.004 & -0.2018 & -0.1255 & -0.1688\\
0.008 & -0.1221 & -0.1436 & -0.2031\\
\bottomrule
\end{tabular}
\end{center}

في القيم الأصغر لـ\(\delta\)، خصوصًا \(0.001\) و\(0.002\)، يتحرك الحد
الثانوي بوضوح نحو البيانات مقارنة بالحد الرئيسي.
ليست هذه مطابقة استدلالية لأن المجال العددي محدود، لكن الإشارة والحجم
متفقان نوعيًا بصورة قوية مع نظرية Leung.

\section{الارتباط عبر مسافات أطول}

تنبؤ Leung هو أن:
\[
\rho_k
\propto
-\Delta(k),
\]
وأن \(\Delta(k)\) تتناقص مع المسافة، وبعيدًا:
\[
\Delta(k)\sim\frac{1}{2k}.
\]

في بياناتنا ظهر بوضوح أن الارتباط الأقوى هو عند الجار، ثم يضعف عند
المسافات \(2\) و\(3\). توجد انحرافات منتهية، خصوصًا عند
\(\delta=0.004,0.008\)، ولذلك لا نستعمل هذه التجربة لإثبات الصيغة؛
بل للتحقق من أننا نقيس الظاهرة نفسها التي تصفها الأدبيات.

\section{اختبار ثلاث فترات دون إدخال بنية ثلاثية}

نظرية Leung لا تقول فقط إن كل زوج له ارتباط؛ بل إن المتجه الكامل
لعدة فترات يصبح تقريبًا Gaussian متعدد المتغيرات.

وهذا مهم لأن Gaussian متعدد المتغيرات لا يحتاج معاملات ثلاثية مستقلة:
بمجرد معرفة مصفوفة التغاير، تتحدد توزيعات الأنماط الثلاثية أيضًا.

اختبرنا ترتيب أول ثلاث فترات.
مجموع حالتي الترتيب الرتيب:
\[
Z_1<Z_2<Z_3
\quad\text{أو}\quad
Z_1>Z_2>Z_3
\]
كان كما يلي:

\begin{center}
\begin{tabular}{rrr}
\toprule
\(\delta\) & المرصود & Gaussian من مصفوفة Leung\\
\midrule
0.001 & 0.3218 & 0.3200\\
0.002 & 0.3119 & 0.3185\\
0.004 & 0.3094 & 0.3162\\
0.008 & 0.3059 & 0.3134\\
\bottomrule
\end{tabular}
\end{center}

كانت أكبر فجوة في احتمال ترتيب منفرد أقل من نحو \(0.012\).

إذن على هذا الاختبار، لا نحتاج إلى معامل ثلاثي مستقل لكي نفسر سلوك
ثلاث فترات؛ مصفوفة الأزواج مع Gaussian approximation تفسر معظم الإشارة.

\section{إعادة فحص بيانات مربعات الأوليات}

عدنا إلى \(450\) فترة:
\[
p_n^2<x<p_{n+1}^2,
\qquad p_n\ge4001.
\]

في تجربة سابقة بنينا مصفوفة تغاير من قانون الفترات القصيرة، وأدخلنا
الارتباطات حتى lag \(50\)، من غير ضبطها على الأنماط الثلاثية المرصودة.

ثم أجرينا \(50\,000\) محاكاة Gaussian.

\subsection{runs}

المرصود:
\[
256.
\]

النموذج الزوجي Gaussian:
\[
240.99\pm10.32.
\]

وقيمة Monte Carlo ثنائية الطرف:
\[
p\approx0.148.
\]

إذن قيمة runs مرتفعة نسبيًا، لكنها ليست شاذة بما يكفي لتكون دليلًا
مقنعًا على ضرورة \(k=3\).

\subsection{الثلاثيات الرتيبة}

من أصل \(448\) ثلاثية متجاورة، كان العدد المرصود للثلاثيات الصاعدة أو
الهابطة رتيبًا:
\[
145.
\]

أما نموذج الأزواج Gaussian فأعطى:
\[
144.55\pm8.86.
\]

والقيمة المرصودة تقع عمليًا في مركز التوزيع المحاكى.

\subsection{أنماط الإشارات الثلاثية}

اختبرنا الأنماط الثمانية:
\[
000,001,010,011,100,101,110,111.
\]

لم يكن أي نمط شاذًا؛ جميع قيم Monte Carlo ثنائية الطرف كانت أكبر من
نحو:
\[
0.22.
\]

إذن نموذجًا يعتمد على covariance زوجية فقط يفسر هذه الإحصاءات الثلاثية
في العينة الحالية بصورة جيدة.

\section{لماذا كان نموذجنا النقطي يوحي بأن triples مطلوبة؟}

هناك فرق بين نموذجين في المشروع:

\begin{enumerate}
\item نموذج count-level Gaussian الذي يأخذ covariance الصحيحة للفترات
القصيرة تقريبًا.
\item نموذج point-level: عجلة ثابتة + random residue tail + thinning.
\end{enumerate}

النموذج الأول أعطى:
\[
\rho_1\approx-0.109
\]
و
\[
\mathrm{runs}\approx241,
\]
وكلاهما منسجم إحصائيًا مع البيانات.

أما أفضل نموذج نقطي في المسح الأخير فبقي تقريبًا عند:
\[
\rho_1\approx-0.083,
\qquad
\mathrm{runs}\approx237.
\]

إذن المشكلة الحالية ليست دليلًا على أن ``الأزواج غير كافية في الطبيعة''؛
بل دليل أقوى على أن \textbf{نموذجنا النقطي لم ينجح بعد في إعادة إنتاج
covariance الزوجية الكاملة التي تعرفها نظرية الفترات القصيرة}.

وهذا يغير ترتيب الخطوات التالية.

\section{موقع Leung 2025}

عمل Leung 2025 عن pseudorandomness يثبت، تحت صيغة موحدة من
Hardy--Littlewood \(k\)-tuple conjecture، تقاربًا إلى Poisson point process
عند مقياس التطبيع بواسطة \(\log N\).

فتراتنا بين مربعات الأوليات أطول بكثير من \(\log N\)، ولذلك التجربة الحالية
أقرب إلى النظام Gaussian للفترات القصيرة من كونها في نظام Poisson المحلي.
هذا يفسر لماذا Leung 2026 وMontgomery--Soundararajan أكثر مباشرة
لمقاييسنا الحالية.

\section{القرار البحثي بعد المقارنة}

القرار السابق كان:
\[
\text{Pair}\longrightarrow\text{Triple}.
\]

بعد المقارنة المباشرة بالأدبيات والاختبار الجديد، القرار الأدق هو:
\[
\boxed{
\text{لا ننتقل إلى triples بعد.}
}
\]

السبب:

\begin{enumerate}
\item الارتباط السلبي المرصود متوافق نوعيًا وكمّيًا مع نظرية Leung للفترات القصيرة.
\item السلوك الثلاثي المحلي يتفق جيدًا مع Gaussian متعدد المتغيرات المحدد
بـpair covariance فقط.
\item حتى runs في بيانات مربعات الأوليات ليس شاذًا تحت نموذج pair-Gaussian.
\item العجز الواضح موجود حاليًا في النموذج النقطي الخاص بالمشروع، الذي لا
يعيد إنتاج قوة covariance الزوجية كاملة.
\end{enumerate}

إذن التجربة التالية الأكثر انضباطًا هي:

\begin{quote}
تعديل النموذج النقطي بحيث يستهدف \emph{مصفوفة covariance الزوجية الكاملة}
عبر عدة lags، لا lag-1 فقط، ثم نعيد اختبار runs والثلاثيات.
\end{quote}

إذا نجح في كل مقاييس الزوج وبقيت إحصاءات ثلاثية شاذة، عندها يصبح الانتقال
إلى Hardy--Littlewood \(k=3\) مبررًا تجريبيًا.

\section{المراجع الأساسية}

\begin{itemize}
\item Sun-Kai Leung,
\emph{Joint distribution of primes in multiple short intervals},
Advances in Mathematics 490 (2026), 110847.
\url{https://www.sciencedirect.com/science/article/pii/S0001870826000691}

\item Sun-Kai Leung,
\emph{Pseudorandomness of primes at large scales},
Quarterly Journal of Mathematics 76 (2025), 251--263.
\url{https://academic.oup.com/qjmath/article/76/1/251/7951830}

\item H. L. Montgomery and K. Soundararajan,
\emph{Primes in short intervals},
Communications in Mathematical Physics 252 (2004), 589--617.
\url{https://arxiv.org/abs/math/0409258}
\end{itemize}

\section{ملفات التجربة}

\begin{itemize}
\item \texttt{theory\_aligned\_comparison\_summary.csv}
\item \texttt{leung2026\_weighted\_interval\_summary.csv}
\item \texttt{leung2026\_lag\_correlation\_comparison.csv}
\item \texttt{leung2026\_triple\_ordering\_comparison.csv}
\item \texttt{prime\_square\_pair\_gaussian\_higher\_order\_diagnostics.csv}
\item \texttt{prime\_square\_pair\_gaussian\_sign\_triples.csv}
\item \texttt{plots/leung2026\_correlation\_comparison.svg}
\item \texttt{plots/leung2026\_runs\_comparison.svg}
\item \texttt{plots/prime\_square\_pair\_gaussian\_sign\_triples.svg}
\end{itemize}

\end{document}
